\section{Yoruba (Ifa)}

The Yoruba tradition is the religion of the Yoruba people of West Africa and of the many communities descended from them across the Atlantic world. It is the most developed and best documented of the African esoteric systems. It holds a remote supreme creator, a wide order of divine powers, a doctrine of personal destiny, and a vast divination system that preserves its teaching as a living oral corpus.

\subsection{The Creation Model}

In the beginning there was only the sky above and a wide marsh of water below, with no solid ground. The supreme creator gave one of the great divine powers the task of making land. That power came down on a chain, carrying a snail shell full of earth, a hen, and a pigeon. He poured out the earth on the water and set the birds down, and they scratched and spread the earth in every direction until it became the wide dry land. The first sacred city was founded on that spot.

The supreme creator then breathed life into the world and set the many divine powers to govern its parts. Each human being, before birth, kneels before the creator and chooses an inner head that carries a destiny. That chosen destiny shapes the life that follows, and the work of a life is to align with it.

\subsection{Key Motifs}

The supreme creator is \textbf{Olodumare}. The power who descends to make land is \textbf{Obatala}. The many divine powers are the \textbf{orisha}, the governors of the forces of the world. The inner head and chosen destiny is \textbf{ori}. The life-force that makes things come to pass is \textbf{ashe}. The teaching is preserved and read through the \textbf{Ifa} divination system, served by the power of wisdom \textbf{Orunmila}. Its corpus is organized into a fixed set of signs.

\begin{longtable}{L{0.22\linewidth}
L{0.6\linewidth}}
\toprule
\textbf{Term} & \textbf{Meaning} \\
\midrule
\textbf{odu} & a sign of the Ifa corpus, each holding many verses \\
\textbf{256 odu} & the full set, formed as sixteen by sixteen \\
\bottomrule
\end{longtable}

This tradition states the shared pattern as the going-out of one remote creator into the many orisha and into the spread of solid land over the water, with the return held in the alignment of each chosen inner head with the destiny set before it.
